\section*{Bab \ref{ch:g3:problems} --- Menyelesaikan Soal Cerita}

\begin{solution}{exo:g3:problems:1}
Kawanan yang bergabung: penjumlahan. Membagi ceri: pembagian. Kotak
berisi pensil: perkalian. Kartu yang diberikan: pengurangan.
\end{solution}

\begin{solution}{exo:g3:problems:2}
Ditanyakan: banyaknya mobil sekarang. Diketahui: $246$, lalu $158$
lagi. Operasinya: penjumlahan. Hitung: $246 + 158 = 404$. Kalimat:
sekarang ada $404$ mobil di tempat parkir itu.
\end{solution}

\begin{solution}{exo:g3:problems:3}
Pembagian pengelompokan: $90 \div 9 = 10$. Talinya memberi $10$ potongan.
\end{solution}

\begin{solution}{exo:g3:problems:4}
Perkalian: $26 \times 8 = 208$. Karcisnya berharga $208$.
\end{solution}

\begin{solution}{exo:g3:problems:5}
Batang pembanding: batang Emi $64$, batang Bayu lebih pendek $29$.
Pengurangan: $64 - 29 = 35$. Bayu punya $35$ stiker.
\end{solution}

\begin{solution}{exo:g3:problems:6}
Yang diterima: $6 \times 45 = 270$ apel. Yang tersisa:
$270 - 178 = 92$ apel.
\end{solution}

\begin{solution}{exo:g3:problems:7}
Buku sains: $6 + 8 = 14$. Novel: $12 - 9 = 3$ lebih banyak pada minggu 2.
\end{solution}

\begin{solution}{exo:g3:problems:8}
Senin $12$, Selasa $20$, Rabu $12$: seluruhnya
$12 + 20 + 12 = 44$ kue. Hari terbaiknya adalah Selasa.
\end{solution}

\begin{solution}{exo:g3:problems:9}
Bilangan yang tidak berguna adalah harganya ($12$). Pembagian:
$210 \div 35 = 6$ (karena $35 \times 6 = 210$). Leo perlu $6$ hari.
\end{solution}

\begin{solution}{exo:g3:problems:10}
Dewasa: $2 \times 9 = 18$. Anak: $3 \times 4 = 12$. Harga seluruhnya:
$18 + 12 = 30$. Kembalian: $50 - 30 = 20$. Mereka menerima $20$ kembali.
\end{solution}

\begin{solution}{exo:g3:problems:11}
Banyak jawaban yang benar --- misalnya: ``Seorang penjual bunga
membuat $8$ ikat berisi $15$ bunga, lalu membuang $20$ bunga yang
layu. Berapa bunga yang tersisa?'' Penyelesaian: $8 \times 15 = 120$,
lalu $120 - 20 = 100$.

\end{solution}
